A arquitetura de microsserviços promete que cada serviço possa mudar e ser implantado sem depender dos demais. Quando vários serviços reutilizam um mesmo componente interno, porém, uma mudança nesse componente pode obrigá-los a mudar juntos. Esse risco é conhecido por relatos de praticantes, mas foi pouco medido em sistemas reais. Este trabalho analisa como o compartilhamento de componentes afeta a autonomia de evolução de microsserviços, com base no código e no histórico de repositórios. A partir de duas bases públicas, com 395 candidatos, oito sistemas Java de código aberto foram sorteados e minerados com o PyDriller. A co-mudança entre serviços, medida pelo índice de Jaccard sobre os \textit{commits}, foi cruzada com as dependências declaradas nos arquivos de \textit{build}, e os 41 componentes compartilhados encontrados foram classificados quanto ao conteúdo. Cerca de metade dos sistemas examinados com três ou mais serviços compartilha código entre eles, e, na mediana, 87\% dos serviços de cada sistema dependem de algum componente. Em cinco dos seis sistemas em que a comparação foi possível, os pares de serviços que compartilham componente mudaram juntos com mais frequência, com diferença significativa e efeito grande em dois. A inspeção de uma amostra de 152 \textit{commits} mostrou, contudo, que só 42\% da co-mudança reflete dependência real, e o restante vem sobretudo de manutenção em lote. A parcela das mudanças de um componente que chega aos serviços não acompanhou o número de serviços que o usam nem a natureza técnica ou de domínio do seu código. As maiores taxas apareceram em módulos pequenos que descrevem a interface de um serviço. Os resultados sugerem que o risco de um componente compartilhado depende menos do seu alcance do que do seu papel como acordo entre serviços.

% Separe as palavras-chave por ponto
\palavraschave{Microsserviços. Reuso de software. Acoplamento lógico. Mineração de repositórios. Componentes compartilhados}
